% 【基础正文】所属：第1章（由 chapters/revised/01-knowledge.tex \input 引入）
\minitopic{JTB分析的吸引力}把知识分析为\term{得到确证的真信念}{justified true belief, JTB}，可以解释三个直觉：知识不能为假，主体必须接受命题，纯粹猜中不算知识。形式化地说：
\[
K_Sp \leftrightarrow p \land B_Sp \land J_Sp .
\]
这里的公式只是分析方案，不是已经证明的定理。尤其“确证”若被定义成“足以产生知识的根据”，分析就会循环。 通常所说的“柏拉图式JTB传统”应谨慎对待。《泰阿泰德》确实讨论真判断加说明，却没有把现代JTB定义确立为此后两千年的无争议定论\parencite[201c--210b]{plato1990}。二十世纪中叶的分析哲学语境，才是Gettier问题的直接背景。

\minitopic{两个原始反例}Gettier只用了三页和两个案例\parencite[121--123]{gettier1963}。第一个是职位与硬币：史密斯从一个得到确证但为假的命题，演绎出一个碰巧为真的命题。第二个案例中，史密斯有强证据相信“琼斯拥有一辆福特车”，并由此任意推出析取命题，如“琼斯拥有福特车，或者布朗在巴塞罗那”。琼斯其实没有福特车，但史密斯毫不知情的布朗恰在巴塞罗那，于是析取命题为真。 两个案例依赖两项假设：一个假命题也可能得到确证；确证在主体知道的演绎推理下传递。Gettier让“有根据”沿推理传下去，而真值从另一条意外路径到达结论。

\begin{argumentbox}
\begin{enumerate}
  \item 若JTB足以构成知识，任何满足真、信念、确证的案例都是知识。
  \item Gettier案例满足这三项条件。
  \item Gettier案例中的主体并不知道目标命题。
  \item 因而JTB不是知识的充分条件。
\end{enumerate}
\end{argumentbox}

结论只是否定充分性。它没有证明真、信念或确证不是必要条件，也没有证明知识不可分析，更没有证明所有确证理论失败。

\minitopic{修补为何困难}最直接的修补是增加“无假前提”：信念不能从任何假命题推得。但有些推理包含无关的假信念，仍可产生知识；反之，一个主体可以不用显式假前提，却被误导性环境包围。例如谷仓立面案例中，亨利看着唯一真谷仓形成“那是谷仓”的真信念；附近其余外观全是假立面。其信念来自真知觉，却仍像是碰巧命中。 “不可击败条件”要求不存在一个真事实，使主体一旦知道它就会失去确证。问题是如何界定真正击败者，避免把“$p$为假”这种平凡信息也算进去。因果理论要求事实以适当方式造成信念\parencite{goldman1967}，能够解释知觉知识，却较难处理数学、道德和普遍命题。 扎格泽布斯基指出，只要理论把获得真理的条件与获得确证的条件分开，就可能人为制造二者意外汇合的案例\parencite{zagzebski1994}。这说明Gettier问题不是寻找第四个词那么简单，而是在问：主体认知成功的\textbf{真}是否源自主体的认识能力或方法。

\minitopic{直觉分歧与方法教训}不同语言和文化中的受试者可能对个别案例作不同判断。由此不能直接推出“知识没有事实”。实验结果需要区分翻译、教育背景、案例复杂度和回答尺度。反过来，哲学家也不应把自己的第一反应当作不可修正的裁决。 Gettier问题的持久意义在于，它迫使理论连接两种评价：信念是否以可辩护方式形成，以及信念为真是否归功于这种方式。后续的安全性、德性和知识优先理论，都是对这项连接任务的不同回答。

\begin{comparison}
“知之为知之”的认识诚实不能单独解决Gettier问题\parencite[2.17]{analects2003}：史密斯并不知道自己处于运气案例，也未必虚假自信。不过，儒家传统把“知”与能否恰当行动、承认无知联系起来，提醒我们：知识概念的实践角色可能超过给命题贴标签。比较应当扩展问题，而不是把经典语句当成现成答案。
\end{comparison}

\minitopic{原典细读：两个假设究竟做了什么}Gettier开篇没有宣称所有确证都可错，而是明确假定：（一）一个得到确证的命题可能为假；（二）若主体对$p$得到确证，且从$p$有效推出$q$并接受$q$，则主体对$q$也得到确证\parencite[121]{gettier1963}。随后两个案例都在这组假设下运作。若把确证定义为事实性的，第一项就会被拒绝；但此时要解释日常所谓“完全有理由却不幸出错”属于什么认识地位。 职位案例中，史密斯并不是推出“我会得到职位”。他推出的是去指称的命题“得到职位的人口袋里有十枚硬币”。这个细节不能省略：原命题以琼斯为根据，结论却可由史密斯意外满足。第二案例使用析取引入规则；史密斯随意选择自己对布朗位置一无所知的地点，正是为了确保结论真值来自完全独立的偶然事实。 准确复述案例是理论训练的一部分。若把结论改成“史密斯会得到职位”，史密斯并没有从原证据得到它，JTB结构就不再成立；若把第二例改成羊或谷仓，则已经换成后来的反例家族，不能再称Gettier的原始案例。

\minitopic{反例谱系}“假谷仓”把问题从错误推理转向环境运气；“停止的钟”中，主体在恰好正确时刻看钟，过程通常不可靠；“假火柴”类案例让主体的证据链含有误导因素，却由另一因素保证真；“鲁莽的认识者”则追问信念虽安全可靠但主体是否缺乏应有责任。不同案例攻击不同修补方案，不能都以“巧合”一句带过。 分析反例可使用四格：（1）目标信念是什么；（2）主体为何相信；（3）何事使它为真；（4）在最近的变化中哪里会出错。若第二项与第三项脱节，通常有Gettier风险；若第四项显示同一方法容易出错，安全性受威胁；若方法稳定但主体忽视明显反证，则责任或信念确证受威胁。

\minitopic{案例工作坊：自动评分}教师使用可靠率99\%的系统评分。某份答案本应得85分，系统因识别错误算成70分；教师又误点修正按钮，恰好把成绩改回85分。学生查看系统后相信自己得85分。信念为真，信息来源通常可靠，但这次真值来自两个错误抵消。因果理论若只要求成绩记录造成信念，可能过宽；一般可靠率若只看长期表现，也会遗漏个案路径；能力或安全理论能更直接指出，该结果不适当地归功于评分过程。 改变细节可检验理论。若教师在发布前逐项人工复核，最终85分便不再是运气；若学生知道系统刚升级且错误频繁，其信念即使为真也可能不理性；若学生只猜自己常得85分，连系统因果链也不存在。通过最小改动观察判断变化，比不断发明离奇故事更能定位知识条件。
